\chapter{实践卷：从 C++ 代码到机器证据}

第一册实践卷围绕：

\begin{lstlisting}[language=bash]
books/cpu-volume-1/labs/lab00_benchmark_foundation
\end{lstlisting}

这一卷训练一条硬链路：C++ 源码、编译选项、可执行文件、输入规模、计时边界、checksum、CSV 报告和反汇编线索。没有这条链路，任何“这个写法更快”的判断都只是感觉。

\section{一、工程入口}

实践工程的核心文件如下：

\begin{lstlisting}[language=bash]
include/lab00/kernels.hpp
include/lab00/benchmark.hpp
src/kernels.cpp
src/benchmark.cpp
src/main.cpp
src/bad_benchmarks.cpp
tests/lab00_tests.cpp
\end{lstlisting}

构建和测试命令如下：

\begin{lstlisting}[language=bash,caption={第一册实践工程构建和测试}]
cmake -S books/cpu-volume-1 -B books/cpu-volume-1/build/lab00-debug -DCMAKE_BUILD_TYPE=Debug
cmake --build books/cpu-volume-1/build/lab00-debug
ctest --test-dir books/cpu-volume-1/build/lab00-debug --output-on-failure
\end{lstlisting}

\section{二、阶段 0：先写 kernel 合同}

\code{kernels.hpp} 中的函数故意很小：

\begin{lstlisting}[language=C++]
std::int64_t sum_i32(std::span<const std::int32_t> values);
float sum_f32(std::span<const float> values);
float dot_f32(std::span<const float> left, std::span<const float> right);
std::uint64_t fill_u8(std::span<std::uint8_t> output, std::uint8_t value);
std::uint64_t copy_u8(std::span<const std::uint8_t> input,
                      std::span<std::uint8_t> output);
\end{lstlisting}

\code{std::span} 是“连续内存的一段视图”：它不拥有数据，只描述起点和长度。这样函数既能接收 \code{std::vector}，也能接收数组切片，而不复制数据。第一阶段只问正确性：求和是否正确，点积长度不一致是否拒绝，填充和复制是否覆盖每个字节。

\section{三、阶段 1：checksum 先守住正确性}

benchmark 最容易犯的错误是“测到的代码其实没做事”。编译器看到某段计算结果没人使用，可能把整段循环删掉。工程里的 \code{checksum_i64}、\code{checksum_f32} 以及每个 workload 返回的 \code{checksum}，就是为了让结果成为可观察事实。

checksum 不是密码学哈希，它在这里是输出摘要。它要满足两个要求：同一输入同一实现反复运行要稳定；不同实现处理同一批输入时，应能用 checksum 检查是否在测同一件事。若 checksum 变化解释不清，时间数据没有资格进入结论。

\section{四、阶段 2：benchmark runner 的边界}

\code{BenchmarkRunner} 的合同包含四个字段：预热次数、测量次数、输入规模和可选的 benchmark 名称过滤。预热不记录结果，测量才写入 \code{BenchmarkResult}。

阶段化实现顺序如下：

\begin{enumerate}
  \item 先实现 \code{parse_size_list}，让 \code{"16, 32,64"} 变成三个正数，并拒绝 0。
  \item 再实现 \code{BenchmarkRunner::add} 和 \code{cases}，确认注册的 case 不丢。
  \item 然后实现 \code{run}，检查 warmup 加 measured 的调用次数，但只记录 measured 结果。
  \item 最后实现 \code{write_csv}，固定表头 \code{name,input_size,iteration,elapsed_ns,checksum}。
\end{enumerate}

测试 \code{test_benchmark_runner_records_measured_iterations_only} 用一个计数 workload 证明 warmup 被执行但不进入结果；\code{test_benchmark_runner_only_filter} 证明过滤条件不会误跑其他 case。

\section{五、阶段 3：正式测量必须说明环境}

Release 构建命令如下：

\begin{lstlisting}[language=bash,caption={运行 Lab 00 benchmark}]
cmake -S books/cpu-volume-1 -B books/cpu-volume-1/build/lab00-release -DCMAKE_BUILD_TYPE=Release
cmake --build books/cpu-volume-1/build/lab00-release --target lab00_bench lab00_bad_benchmarks lab00_tests
ctest --test-dir books/cpu-volume-1/build/lab00-release --output-on-failure
books/cpu-volume-1/build/lab00-release/bin/cpu-volume-1/lab00/lab00_bench --sizes 1024,65536,1048576
\end{lstlisting}

报告至少写清 CPU 型号、编译器版本、构建类型、输入规模、预热次数、测量次数、是否有后台负载、是否固定频率。Debug 和 Release 的机器码可能完全不同；虚拟机、远端机器和手机环境也可能因为调度、温度或频率策略带来噪声。

\section{六、阶段 4：从源码回到汇编}

第一册正文讲 C++、ABI、汇编和机器执行，实践卷就不能只停在 CSV。每个性能结论至少要回看一个机器证据：

\begin{enumerate}
  \item 热循环里有没有意外函数调用；
  \item 是否出现向量指令；
  \item load/store 是否连续；
  \item 分支是否在内层循环里；
  \item 编译器是否把计算删除或合并。
\end{enumerate}

可以用 \cmd{objdump -d}、\cmd{perf} 或编译器生成汇编。若平台没有工具，报告要写“未验证”，不能把推断写成事实。

\section{七、阶段 5：坏 benchmark 也要入书}

\code{src/bad_benchmarks.cpp} 的作用是教你识别坏测量。典型坏例子是在计时循环里反复分配大向量，结果仍然可能有正确 checksum，但测到的主要是分配器、缺页、缓存扰动和初始化成本，不是 \code{dot_f32} 本身。

写坏 benchmark 报告时分三步：先证明它的输出语义没有错；再说明它把哪些额外工作放进了计时边界；最后给出修正版，把输入准备移到计时循环之外。这样读者不会把“代码能跑”误解为“测量可信”。

\section{八、项目阶段清单}

\begin{longtable}{p{0.16\linewidth}p{0.28\linewidth}p{0.46\linewidth}}
阶段 & 代码切片 & 验收证据 \\
\hline
0 & \code{sum_i32/sum_f32/dot_f32/fill_u8/copy_u8} & 单元测试检查数值、长度错误和字节结果 \\
1 & \code{checksum_i64/checksum_f32} & benchmark workload 返回稳定摘要 \\
2 & \code{BenchmarkRunner} & warmup 不入结果，measured 写入结果 \\
3 & \code{write_csv} 与命令行 sizes & CSV 表头稳定，非法 sizes 被拒绝 \\
4 & Release benchmark & 报告记录环境、输入、迭代和 checksum \\
5 & 反汇编或 perf 回看 & 至少解释一个热循环机器证据 \\
6 & \code{lab00_bad_benchmarks} & 说明坏测量把哪些无关成本放进边界 \\
\end{longtable}

\begin{labbox}
扩展任务：给 \code{dot_f32} 增加一个“错误 benchmark”版本，例如在测量循环中重复分配向量。先证明它也能得到正确 checksum，再解释为什么它测到的是分配器和缓存扰动，不是纯 dot product。最后写出修正版，并把两者 CSV 放在报告中对比。
\end{labbox}
